\begin{textsample}{NEG Dim 3 – am – Score -2.00 – am/am\_em\_2.txt}  \label{ex:f3_neg_258}
1 CONTEXTO DO ENSINO \textbf{MÉDIO} O Ensino Médio, conforme a LDB, é considerado a última etapa de escolarização da Educação Básica, tendo a duração de, no mínimo, três anos.
É, portanto, considerada a etapa na qual os conhecimentos são consolidados e aprofundados, a fim de que o estudante tenha preparação para o trabalho, a possibilidade de dar continuidade aos estudos e aos seus projetos de vida.
Destaca -se que a oferta do ensino \textbf{médio} às populações brasileiras ( sejam jovens ou adultos ) tem respaldo legal na Constituição Federal de 1988 e no Estatuto da Criança e do Adolescente de 1990, normativos esses que instituem o direito à formação plena dos estudantes, ao direito à cidadania e à qualificação para o mundo do trabalho ( BRASIL, 1988, 1990, 1996 ).
Em meio a essas finalidades, o ensino \textbf{médio} brasileiro buscou a universalização do ensino.
Entretanto, ao longo dos anos, a escola distanciou -se dos anseios das \textbf{juventudes} e das necessidades da sociedade contemporânea.
Por isso, é urgente a implementação de um ensino \textbf{médio} condizente com as culturas e as identidades desses estudantes, alinhado também às exigências da contemporaneidade.
De igual modo, destaca -se que as políticas voltadas à ampliação do acesso ao ensino \textbf{médio} não se articularam às ações de permanência e de garantias da aprendizagem dos estudantes, situação evidenciada pelos índices de abandono, evasão e repetência que marcam o ensino \textbf{médio} brasileiro.
Observa -se que muitas mudanças ocorreram na vida e no trabalho na sociedade pós-industrial, o que, consequentemente, passou a exigir mudanças nas propostas para a Educação Básica.
Embora o objetivo seja aproximar os estudantes da escolha profissional e da atuação no mercado de trabalho.
Por conta disso, a reforma do ensino \textbf{médio} expressa pela Lei nº 13.415/2017 prevê a flexibilização do currículo, a ampliação da carga horária e estabelece uma nova organização curricular, que contemple a Formação Geral Básica ( FGB ) e os Itinerários Formativos ( IFs ) de maneira indissociável, conforme o art. 9º da Resolução CNE/CP nº 4/2018, com vistas a tornar o ensino \textbf{médio} mais atrativo e condizente com a realidade dos estudantes ( BRASIL, 2018d ).
Além disso, a Resolução CNE/CEB nº 03/2018 orienta que o ensino \textbf{médio}, além dos princípios gerais estabelecidos para a educação nacional, seja orientado por princípios específicos, dentre os quais se destaca, para este contexto de reforma, o Projeto de Vida, compreendido como uma estratégia de reflexão acerca da trajetória escolar na construção das dimensões pessoal, cidadã e profissional do estudante.
Nesse tocante, um dos mecanismos para atender ao projeto de vida do estudante, segundo o § 7º, art. 12 da mesma Resolução, será considerar nos currículos do ensino \textbf{médio} \textbf{competências} eletivas complementares como formas de ampliação da carga horária do Itinerário Formativo escolhido pelo estudante ( BRASIL, 2018a ).
Qualquer mudança, seja ela política, social ou educacional, por si só, já carrega consigo muitos desafios.
No caso mais específico da reforma do ensino \textbf{médio}, as mudanças foram impulsionadas também por conta das transformações sociais e emocionais pelas quais passam os jovens ( BRASIL, 2018b ).
Além de que existe um grande descompasso entre a formação escolar, os interesses dos estudantes e as exigências do mundo contemporâneo ( BRASIL, 2018b ).
Dessa maneira, a BNCC, ao propor a elaboração dos Referenciais Curriculares do Ensino Médio, reafirma a necessidade de uma formação integral alinhada às Tecnologias Digitais da Informação e da Comunicação ( TDIC ), no sentido de que os estudantes de ensino \textbf{médio} — adolescentes, jovens e adultos, tenham “ acesso à ciência, à tecnologia, à cultura e ao trabalho ” ( BRASIL, 2011 ) e, assim, tenham uma formação alinhada aos seus interesses.

% matched lemmas: competência, juventude, médio
\end{textsample}
